\chapter{Estudo das tecnoloxías para a creación de aplicacións web baseadas en vídeo}

	A continuación mostrase un estudo feito sobre as tecnoloxías mais empregadas hoxe en día
	de cara ao deseño de aplicacións web baseadas en vídeo. Os criterios que se valoraron á
	hora de escoller unha destas metodoloxías son os seguintes:

\begin{itemize}

	\item {\textbf{Plataforma e Portabilidade\\}}
		Segundo as especificacións iniciais do proxecto, e de cara a facilitar a 
		implementación deste, as ferramentas empregadas deben de executarse baixo os
		Sistemas Operativos baseados en Linux.

	\item {\textbf{Compatibilidade co algoritmo de Procesado de Vídeo\\}}
		Dado que o algoritmo que procesa o vídeo foi parcialmente implementado en c++,
		a tecnoloxía que se empregue para o desenvolvemento da parte web debe dispoñer da
		máxima compatibilidade con esta linguaxe de programación.

	\item {\textbf{Desenvolvemento Áxil\\}}
		A tecnoloxía que se seleccione ten que minimizar o custe en tempo e esforzo da 
		implementación, por este motivo valorarase positivamente que dispoña de IDE's axeitados,
		facilidades de acceso a BD(Base de Datos),se é posible ORM(Object-Relational Mapping) 
		integrado, recarga en quente... En xeral todo aquelo que permita axilidade e flexibilidade.

\end{itemize}

\section{Back-End}
	O Back-End é a parte que procesa os datos proporcionados polo Front-End, atende as súas 
	peticións, e xestiona o modelo de datos e os procesos implicados na aplicación. Neste
	caso en concreto, o Back-End será o encargado de interactuar directamente co sistema de
	Análise de Vídeo.\\
	
	As tecnoloxías estudadas para esta parte do sistema son as seguintes:

	\subsection{Java}
  	Java é unha das linguaxes mais empregadas actualmente. Ademais existen diversos
  	frameworks web como Tapestry ou SpringMVC e en canto as BD Hibernate, que facilitan o 
  	seu uso, mais é preciso integralos xa que non forman parte da plataforma en si. A súas 
  	posibilidades de integración con c++ son altas grazas á interface JNI(Java Native 
  	Interface), pero a súa configuración pódese volver tediosa. Un dos proxectos estudados
  	para ver o seu funcionamento é Red5 \cite{red5-github-url}.
  	
	\subsection{C\#}
	esta linguaxe en combinación con .NET resulta unha combinación bastante áxil de cara
	á programación web, integrando na propia plataforma un deseñador Web e un ORM moi 
	intuitivo. O gran problema polo que se descartou este entorno foi polo seu baixo grao
	de compatibilidade cos sistemas operativos Linux.
	
	\subsection{C - C++}
	C++ presenta como era de esperar a maior compatibilidade co algoritmo implementado, non
	obstante, inda que existen algunhas utilidades que facilitan o desenvolvemento web con 
	esta linguaxe como Wt (Web Toolkit)\cite{wt-url}, o grado de axilidade está moi por 
	baixo do que facilitan o resto das combinacións. Tamén pode resultar de interese o
	coñecido proxecto Icecast\cite{icecast-url}, que fai streaming de vídeo sobre unha 
	interface web.
	
	\subsection{Python + Django}
	Python presentase como a mellor opción para desenvolver o lado servidor, por unha parte
	dispón do módulo Subprocess\cite{subprocess-module-url} que permite executar calquera
	comando pola terminal maximizando así a modularidade e a integración co algoritmo en C++.
	Por outra parte Django\cite{django-web-page-url} contén un potente ORM e un sistema de 
	"Templates" que simplifica a parte web. Para concluír cabe destacar o feito de que 
	python sexa unha linguaxe interpretada, xa que isto evita o paso previo de compilación.
	
	
\section{Front-End}
	É a parte do Sistema encargada de interactuar co usuario. O maior requisito de esta 
	parte é ser capaz de representar todas as capas de información complexa que xera o
	Back-End. E tratándose dunha aplicación web estudaranse aqueles xeitos que permitan
	a reprodución de vídeo, e a maiores o debuxado de figuras en movemento sobre este vídeo.
	
	\subsection{Vídeo en Flash}
		Video Flash é a tecnoloxía de reprodución de vídeo mais empregada e madura en
		internet dende hai anos. Inicialmente creada por Macromedia e mercada por Adobe 
		en 2005, permite crear elaboradas animacións vectoriais, que logo poden proxectarse
		sobre un vídeo, e á vez tamén manexa os eventos de reprodución de vídeo como o 
		play ou o stop.\\
		
		Ten certos problemas en tanto ao Posicionamento Web(SEO), reprodución en 
		dispositivos móbiles, accesibilidade... pero o meirande de todos eles é que mentres
		que outros dos exemplos estudados son 100\% libres, Flash é un programa propietario
		para o que é preciso adquirir unha licencia.		
		
	
	\subsection{Vídeo HTML5 + Javascript}
		Está é unha das combinacións mais empregadas actualmente, xa que segue o estándar
		do W3C (World Wide Web Consortium)\cite{w3schools-video-tag}. Nel defínese como se
		han de mostrar e obter os vídeos dunha páxina codificada coa linguaxe HTML5, e 
		destaca por seres extremadamente sinxelo en comparación con Flash ou outras tecnoloxías.
		
		Para a modificación da súa aparencia pódese empregar \emph{CSS3} (Follas de Estilos en 
		Cascada), e para modificar o seu comportamento, o elemento $<video>$ de HTML5 achega
		unha serie de Métodos, Eventos e Propiedades\cite{w3school-video-events} que poden
		empregarse dende código \emph{Javascript}.
		
		Isto resulta de moita utilidade, xa que mediante Javascript tamén pode manexarse o
		elemento $<canvas>$ de HTML5, que xera un mapa de bits para construír gráficos, 
		manipular imaxes e crear dinamicamente animacións nunha páxina web. A única dúbida que
		soe xurdir sobre esta tecnoloxía está en canto ao seu rendemento e alcance, pero
		exemplos como os que amosa Kevin Roast na súa páxina web\cite{kevin-roast-canvas-examples}
		despexan toda dúbida posíbel. 
		
		Tamén é de interese a librería de estilos web \emph{Twiter Bootstrap}
		\cite{bootstrap-page-url} para axilizar o deseño web con CSS3, Así como a 
		librería de Javascript \emph{JQuery} \cite{jquery-page-url} que minimiza de 
		xeito considerable as liñas de código a implementar.
		

\vspace{2cm}
%%%%%%%%%%%%%%%%%%%%%%%%%%%%%%%%%%%%%%%%%%%%%%%%%%%%%%%%%%%%%%%%%%%%%%%%%%%%%%%%%%

Esta es una referencia a un artículo \cite{article-minimal}.

Esta es una referencia a un capítulo dentro de un libro \cite{inbook-minimal}.

Esta es una referencia a un libro \cite{book-minimal}.

Esta es una referencia a un artículo dentro de los proceedings de un congreso \cite{inproceedings-full}.

Esta es una referencia a una url \cite{misc-url}.

\textbf{Esto está escrito en negrita}


\emph{Esto está en enfatizado}

\begin{center}
Este texto está centrado
\end{center}


Esto es una lista:
\begin{itemize}
\item Primer elemento
\item Segundo elemento
\item Lista dentro de otra lista:
	\begin{itemize}
		\item Primer subelemento
		\item Segundo subelemento
	\end{itemize}
\end{itemize}


\section{Título sección}

Esto es una descripción:
\begin{description}
\item[Palabra] descripción de la palabra
\item[Palabra] descripción de la palabra
\end{description}


Y esto una lista numerada:
\begin{enumerate}
\item Elemento
\item Elemento
\item Elemento
\item Elemento
\end{enumerate}


Podemos incluir una figura y referenciarla de esta forma \ref{fig:logo}. Además podemos poner la página en la que está: \pageref{fig:logo}

\begin{figure}[htp]
\begin{center}
\includegraphics[scale=0.2]{figures/logo.eps}
\caption{Esta es la etiqueta de la figura}
\label{fig:logo}
\end{center}
\end{figure}


\begin{figure}[htp]
\begin{center}
\framebox{\includegraphics[scale=0.2]{figures/logo.eps}}
\caption{Esta es la etiqueta de la figura con borde}
\label{fig:logo2}
\end{center}
\end{figure}


\subsection{Título subsección}

Existen varias formas de incluir ecuaciones matemáticas. Las más utilizadas son las siguientes:
\begin{itemize}
 \item Para introducir expresiones matemáticas en el texto, se utiliza como delimitador el símbolo del dólar. Por ejemplo $a \rightarrow b$.
 \item Para introducir ecuaciones matemáticas, se utiliza el entorno \texttt{equation}:
\begin{equation}
 \gamma = \frac{\overline{\alpha}}{\sqrt{\beta}}
\label{eq:equation_example_1}
\end{equation}

\begin{equation}
E(v) =  \int^1_0 \int^1_0 \int^1_0 E_{int}(v(r,s,t)) + E_{ext}(v(r,s,t))drdsdt 
\label{eq:equation_example_2}
\end{equation}	
Un subíndice se especifica con el guión bajo y un superíndice con el circunflejo. Si el super/subíndice contiene varios caracteres, estos deben estar delimitados por llaves. Consultar el manual para comprobar como se pueden introducir símbolos y expresiones matemáticas en latex.
\end{itemize}



\subsection{Título subsección}

Esto es una subsección\footnote{Así se hace una nota a pie de página}.


Así introducimos texto sin ningún tipo de formato latex:
\begin{verbatim}
  4 drwxr-xr-x  2 noelia imagen   4096 2005-09-12 12:09 figures
  4 -rwxr--r--  1 noelia imagen    585 2005-09-12 16:56 Makefile
  4 -rw-r--r--  1 noelia imagen    647 2005-09-12 17:38 memoria.aux
  4 -rw-r--r--  1 noelia imagen   1011 2005-09-12 17:18 memoria.bbl
  4 -rw-r--r--  1 noelia imagen   1171 2005-09-12 17:18 memoria.blg
 16 -rw-r--r--  1 noelia imagen  13440 2005-09-12 17:38 memoria.dvi
  4 -rw-r--r--  1 noelia imagen    412 2005-09-12 17:38 memoria.glg
  4 -rw-r--r--  1 noelia imagen    188 2005-09-12 17:38 memoria.glo
  4 -rw-r--r--  1 noelia imagen    241 2005-09-12 17:38 memoria.gls
  4 -rw-r--r--  1 noelia imagen    299 2005-09-12 17:38 memoria.ist
  4 -rw-r--r--  1 noelia imagen    283 2005-09-12 17:38 memoria.lof
 12 -rw-r--r--  1 noelia imagen   8997 2005-09-12 17:38 memoria.log
  4 -rw-r--r--  1 noelia imagen     87 2005-09-12 17:38 memoria.lot
108 -rw-r--r--  1 noelia imagen 103492 2005-09-12 17:18 memoria.pdf
280 -rw-r--r--  1 noelia imagen 278687 2005-09-12 17:18 memoria.ps
  4 -rwxr--r--  1 noelia imagen   1887 2005-09-12 17:40 memoria.tex
  4 -rw-r--r--  1 noelia imagen    735 2005-09-12 17:38 memoria.toc
\end{verbatim}



\subsubsection{Título subsubsección}

La figura \ref{tab:tabla_ejemplo} muestra un ejemplo de tabla básica. Para crear una tabla se utiliza el entorno \texttt{tabular} dentro del flotante \texttt{table}. 

%% PARA ENTENDER ESTE TEXTO, MEJOR COMPILAR Y LEER EL DVI/PS/PDF
%% \ es un caracter reservado de latex por lo que para poder utilizarlo cuando escribimos debemos utilizar el comando \textbackslash
%% &, {, } también son caracteres reservados. En este caso, para escribirlos en el texto, anteponemos \ al caracter

Tras \texttt{\textbackslash begin\{tabular\}} se declara el número de columnas de la tabla. Cada columna se especifica con una letra:
\begin{description}
 \item [c] Si el texto en la columna está centrado
 \item [l] Si el texto en la columna está alineado a la izquierda
 \item [r] Si el texto en la columna está alineado a la derecha
\end{description}

\begin{table}[htp]
\caption{Tabla de ejemplo}
\label{tab:tabla_ejemplo}
\begin{center}
 \begin{tabular}{rlc}\hline
  Fila 1 Col. 1 larálala & Fila 1 Col. 2 lara & Fila 1 Col. 3  lalalalalalalala\\ \hline
  Fila 2 Col. 1 & Fila 2 Col. 2 & Fila 2 Col. 3 \\ \hline
 \end{tabular}
\end{center}
\end{table}

Así \texttt{\textbackslash begin\{tabular\}\{rrrr\}} indica 4 columnas con texto alineado a la derecha y \texttt{\textbackslash begin\{tabular\}\{cl\}} indica dos columnas, la primera centrada y la segunda alineada a la izquierda.

Entre \texttt{\textbackslash begin\{tabular\}\{\ldots\}} y \texttt{\textbackslash end\{tabular\}} se escribe el contenido de la tabla. Las columnas se separan con \& y las filas con \textbackslash \textbackslash. Un ejemplo de una fila con tres columnas sería el siguiente:
\begin{center}
 \texttt{ aaa \& bbbb \& cccc \textbackslash \textbackslash}\footnote{Importante: debe coincidir el número de columnas en la declaración con el número de columnas que se escriben dentro de la tabla}. 
\end{center}

Para incluir lineas horizontales existe el comando \texttt{\textbackslash hline}. Tras \texttt{\textbackslash begin\{tabular\}\{\ldots\}} podemos incluir un \texttt{\textbackslash hline}, pero \textbf{ojo}, en cada fila \texttt{\textbackslash hline} siempre debe ir tras \textbackslash \textbackslash.

También es posible incluir líneas verticales, pero los manuales de estilo no lo aconsejan. Las líneas verticales se definen en la declaración con barras verticales. Por ejemplo, con \{\textbar c  \textbar c \textbar c\}  se crearían 3 líneas verticales, una antes de la primera columna, otra entre la primera y la segunda columna, y la tercera, entre la segunda y tercera columna.

Dentro de las tablas podemos incluir expresiones matemáticas, texto enfatizado, negrita, etc. Es posible también fusionar varias filas o varias columnas dentro de una misma tabla. Para ello, consultar en manuales los comandos/paquetes \texttt{multirow} y \texttt{multicol}.


